% GENERATED FILE. Edit canonical lessons, then run npm run generate:books.
% canonical-source-hash: fnv1a64:aafb02786b8cdb38
% canonical-lessons: JA-C288-nyuugaku, JA-C288-shiki, JA-C288-shusseki-suru, JA-C288-nigiyaka, JA-C288-oiwai

\chapter{Starting School}
\label{ch:ja-starting-school}
\hlchaptermodality{288}

\begin{chapteropening}
\textbf{By the end of this chapter:} \emph{I can say who started school and who came to the ceremony.}

Complete the last lesson of chapter 288: I can say who started school and who came to the ceremony.
\end{chapteropening}

\section[\texorpdfstring{\ja{にゅうがく} (nyūgaku)}{にゅうがく (nyūgaku)}]{\ja{にゅうがく} (nyūgaku) — entering school}

\label{lesson:JA-C288-nyuugaku}



\begin{quote}
Before the new one: say the Japanese for a relative, then the Japanese for to grow up.
\end{quote}

\subsection*{You'll want to know: \ja{にゅうがく}}
\textbf{\ja{にゅうがく}} — \emph{nyūgaku} — \textquotedblleft{}entering school\textquotedblright{}.

\begin{quote}
\textbf{\ja{おとうとは} \ja{らいねん} \ja{にゅうがくします。}} — \emph{otōto wa rainen nyūgaku shimasu} — \textquotedblleft{}My little brother starts school next year.\textquotedblright{}
\end{quote}

\begin{cousinweb}
\textbf{\ja{にゅうがく}} (\emph{nyūgaku}) is Sino-Japanese (kango): two Chinese characters read with their on'yomi, readings borrowed from Chinese. They are \emph{nyū}, \textquotedblleft{}enter\textquotedblright{} (Middle Chinese \emph{nyip}) and \emph{gaku}, \textquotedblleft{}learning\textquotedblright{} (\emph{haewk}) (KANJIDIC2; Middle Chinese in Baxter's transcription).
\end{cousinweb}

\begin{grammarlens}[title={What you've built}]
One more word for telling what happened.
\end{grammarlens}

\subsection*{Your turn}
\begin{itemize}
\raggedright
  \item \emph{Say it:} \emph{nyūgaku}
  \item \emph{Say it:} \emph{nyūgaku}, once more
  \item \emph{Say it:} \emph{otōto wa rainen nyūgaku shimasu} — My little brother starts school next year.
  \item \emph{Recall:} say the Japanese for a relative, then the Japanese for to grow up, then say \emph{nyūgaku} again
\end{itemize}

\begin{tcolorbox}[breakable,colback=teal!4,colframe=teal!35!black,title={Before you move on}]
What does \ja{にゅうがく} mean? (\textquotedblleft{}Entering school\textquotedblright{}.) Say it once more.
\end{tcolorbox}
\section[\texorpdfstring{\ja{しき} (shiki)}{しき (shiki)}]{\ja{しき} (shiki) — a ceremony}

\label{lesson:JA-C288-shiki}



\begin{quote}
Before the new one: say the Japanese for to grow up, then the Japanese for entering school.
\end{quote}

\subsection*{You'll want to know: \ja{しき}}
\textbf{\ja{しき}} — \emph{shiki} — \textquotedblleft{}a ceremony\textquotedblright{}.

\begin{quote}
\textbf{\ja{がっこうで} \ja{しきが} \ja{ありました。}} — \emph{gakkō de shiki ga arimashita} — \textquotedblleft{}There was a ceremony at the school.\textquotedblright{}
\end{quote}

\begin{cousinweb}
\textbf{\ja{しき}} (\emph{shiki}) is Sino-Japanese (kango): one Chinese character read with its on'yomi, a reading borrowed from Chinese. It is \emph{shiki}, \textquotedblleft{}ceremony\textquotedblright{} (Middle Chinese \emph{syik}) (KANJIDIC2; Middle Chinese in Baxter's transcription).
\end{cousinweb}

\begin{grammarlens}[title={What you've built}]
One more word for telling what happened.
\end{grammarlens}

\subsection*{Your turn}
\begin{itemize}
\raggedright
  \item \emph{Say it:} \emph{shiki}
  \item \emph{Say it:} \emph{shiki}, once more
  \item \emph{Say it:} \emph{gakkō de shiki ga arimashita} — There was a ceremony at the school.
  \item \emph{Recall:} say the Japanese for to grow up, then the Japanese for entering school, then say \emph{shiki} again
\end{itemize}

\begin{tcolorbox}[breakable,colback=teal!4,colframe=teal!35!black,title={Before you move on}]
What does \ja{しき} mean? (\textquotedblleft{}A ceremony\textquotedblright{}.) Say it once more.
\end{tcolorbox}
\section[\texorpdfstring{\ja{しゅっせきする} (shusseki suru)}{しゅっせきする (shusseki suru)}]{\ja{しゅっせきする} (shusseki suru) — to attend; to be present}

\label{lesson:JA-C288-shusseki-suru}



\begin{quote}
Before the new one: say the Japanese for entering school, then the Japanese for a ceremony.
\end{quote}

\subsection*{You'll want to know: \ja{しゅっせきする}}
\textbf{\ja{しゅっせきする}} — \emph{shusseki suru} — \textquotedblleft{}to attend; to be present\textquotedblright{}.

\begin{quote}
\textbf{\ja{ははも} \ja{しきに} \ja{しゅっせきします。}} — \emph{haha mo shiki ni shusseki shimasu} — \textquotedblleft{}My mother is going to the ceremony too.\textquotedblright{}
\end{quote}

Polite \textbf{\ja{しゅっせきします}}, \emph{shussekishimasu}; a \ja{する} verb. Past (polite): \textbf{\ja{ちちが} \ja{しきに} \ja{しゅっせきしました。}}, \emph{chichi ga shiki ni shusseki shimashita}.

\begin{cousinweb}
\textbf{\ja{しゅっせきする}} (\emph{shusseki suru}) is Sino-Japanese (kango): two Chinese characters read with their on'yomi, readings borrowed from Chinese. They are \emph{shutsu}, \textquotedblleft{}exit\textquotedblright{} (Middle Chinese \emph{tsyhwit}) and \emph{seki}, \textquotedblleft{}seat\textquotedblright{} (\emph{zjek}) (KANJIDIC2; Middle Chinese in Baxter's transcription). The first is shortened from \emph{shutsu} to \emph{shus-} before \emph{seki}. The native verb \emph{suru}, \textquotedblleft{}to do\textquotedblright{}, makes it a verb.
\end{cousinweb}

\begin{grammarlens}[title={What you've built}]
One more word for telling what happened.
\end{grammarlens}

\subsection*{Your turn}
\begin{itemize}
\raggedright
  \item \emph{Say it:} \emph{shusseki suru}
  \item \emph{Say it:} \emph{shusseki suru}, once more
  \item \emph{Say it:} \emph{haha mo shiki ni shusseki shimasu} — My mother is going to the ceremony too.
  \item \emph{Recall:} say the Japanese for entering school, then the Japanese for a ceremony, then say \emph{shusseki suru} again
\end{itemize}

\begin{tcolorbox}[breakable,colback=teal!4,colframe=teal!35!black,title={Before you move on}]
What does \ja{しゅっせきする} mean? (\textquotedblleft{}To attend; to be present\textquotedblright{}.) Say it once more.
\end{tcolorbox}
\section[\texorpdfstring{\ja{にぎやか} (nigiyaka)}{にぎやか (nigiyaka)}]{\ja{にぎやか} (nigiyaka) — lively; bustling}

\label{lesson:JA-C288-nigiyaka}



\begin{quote}
Before the new one: say the Japanese for a ceremony, then the Japanese for to attend.
\end{quote}

\subsection*{You'll want to know: \ja{にぎやか}}
\textbf{\ja{にぎやか}} — \emph{nigiyaka} — \textquotedblleft{}lively; bustling\textquotedblright{}.

\begin{quote}
\textbf{\ja{この} \ja{まちは} \ja{年}\ja{じゅう} \ja{にぎやかです。}} — \emph{kono machi wa nenjū nigiyaka desu} — \textquotedblleft{}This town is lively all year round.\textquotedblright{}
\end{quote}

A na-adjective: \ja{な} before a noun.

\begin{cousinweb}
\textbf{\ja{にぎやか}} (\emph{nigiyaka}) is a native Japanese word (wago), inherited rather than borrowed from Chinese; its kanji was attached later as a kun'yomi (JMdict).
\end{cousinweb}

\begin{grammarlens}[title={What you've built}]
One more word for telling what happened.
\end{grammarlens}

\subsection*{Your turn}
\begin{itemize}
\raggedright
  \item \emph{Say it:} \emph{nigiyaka}
  \item \emph{Say it:} \emph{nigiyaka}, once more
  \item \emph{Say it:} \emph{kono machi wa nenjū nigiyaka desu} — This town is lively all year round.
  \item \emph{Recall:} say the Japanese for a ceremony, then the Japanese for to attend, then say \emph{nigiyaka} again
\end{itemize}

\begin{tcolorbox}[breakable,colback=teal!4,colframe=teal!35!black,title={Before you move on}]
What does \ja{にぎやか} mean? (\textquotedblleft{}Lively; bustling\textquotedblright{}.) Say it once more.
\end{tcolorbox}
\section[\texorpdfstring{\ja{おいわい} (oiwai)}{おいわい (oiwai)}]{\ja{おいわい} (oiwai) — a celebration; a present for a happy occasion}

\label{lesson:JA-C288-oiwai}



\begin{quote}
Before the new one: say the Japanese for to attend, then the Japanese for lively.
\end{quote}

\subsection*{You'll want to know: \ja{おいわい}}
\textbf{\ja{おいわい}} — \emph{oiwai} — \textquotedblleft{}a celebration; a present for a happy occasion\textquotedblright{}.

\begin{quote}
\textbf{\ja{おいわいに} \ja{みんなで} \ja{しょくじを} \ja{しました。}} — \emph{oiwai ni minna de shokuji o shimashita} — \textquotedblleft{}We all had a meal to celebrate.\textquotedblright{}
\end{quote}

\begin{cousinweb}
\textbf{\ja{おいわい}} (\emph{oiwai}) is a native Japanese word (wago), inherited rather than borrowed from Chinese, built from \emph{o-}, \textquotedblleft{}the polite prefix\textquotedblright{}, and \emph{iwai}, \textquotedblleft{}celebrating\textquotedblright{} (from \emph{iwau}) (JMdict, KANJIDIC2).
\end{cousinweb}

\begin{grammarlens}[title={What you've built}]
One more word for telling what happened.
\end{grammarlens}

\subsection*{Your turn}
\begin{itemize}
\raggedright
  \item \emph{Say it:} \emph{oiwai}
  \item \emph{Say it:} \emph{oiwai}, once more
  \item \emph{Say it:} \emph{oiwai ni minna de shokuji o shimashita} — We all had a meal to celebrate.
  \item \emph{Recall:} say the Japanese for to attend, then the Japanese for lively, then say \emph{oiwai} again
\end{itemize}

\begin{tcolorbox}[breakable,colback=teal!4,colframe=teal!35!black,title={Before you move on}]
What does \ja{おいわい} mean? (\textquotedblleft{}A celebration; a present for a happy occasion\textquotedblright{}.) Say it once more.
\end{tcolorbox}
